% Generated by scripts/assemble.py; corrections belong in source/corrections.json.
\examhead{49}{东华大学 2026 年数学分析试卷}{50}
% p051-b19; source PDF page 51
\par\Needspace{5\baselineskip}
\blockanchor{p051-b19}
\majorquestion{一}计算题. 每题8分, 共40分.\par
% p051-b20; source PDF page 51
\blockanchor{p051-b20}
\question{1}求极限 $\lim_{n\to \infty}\sum_{i = 1}^{n}\sin \left(\frac{i}{n^2}\right)\ln \left(1 + \frac{i}{n}\right).$\par
% p051-b21; source PDF page 51
\blockanchor{p051-b21}
\question{2}求极限 $\lim_{x\to 0}\frac{1}{x}\left(\frac{1}{x} -\cot x\right).$\par
% p051-b22; source PDF page 51
\blockanchor{p051-b22}
\question{3}求积分 $\int_0^1\frac{\ln x}{1 - x}\mathrm{d}x.$\par
% p051-b23; source PDF page 51
\blockanchor{p051-b23}
\question{4}求幂级数 $\sum_{n=0}^{\infty} \frac{x^{2n+2}}{(n+1)(2n+1)}$ 的收敛域及和函数.\par
% p051-b24; source PDF page 51
\blockanchor{p051-b24}
\question{5}求积分 $\iint_{\Sigma} \frac{x \, \mathrm{d}y \, \mathrm{d}z + y \, \mathrm{d}z \, \mathrm{d}x + z \, \mathrm{d}x \, \mathrm{d}y}{(x^2 + y^2 + z^2)^{\frac{3}{2}}}$, 其中 $\Sigma$ 为椭球面 $x^2 + 2y^2 + 3z^2 = 6$ 的内侧.\par
% p052-b00; source PDF page 52
\par\Needspace{5\baselineskip}
\blockanchor{p052-b00}
\majorquestion{二}（15分）设实数列 $\{a_{n}\}$ 满足 $a_1 > 0, a_{n+1} = a_n + \frac{1}{na_n}, n = 1,2,\cdots$. 证明：$\lim_{n\to \infty}\frac{a_n^2}{\ln n} = 2.$\par
% p052-b01; source PDF page 52
\par\Needspace{5\baselineskip}
\blockanchor{p052-b01}
\majorquestion{三}(20 分) 叙述实数完备性的有限覆盖定理, 并用之证明一元函数的介值性定理.\par
% p052-b02; source PDF page 52
\par\Needspace{5\baselineskip}
\blockanchor{p052-b02}
\majorquestion{四}（15分）设函数 $f(x)$ 在 $(- \infty, + \infty)$ 上具有二阶连续导数，证明：$f''(x) \geq 0$ 的充要条件是对任意不同实数 $a, b,$ 有 $f\left(\frac{a + b}{2}\right) \leq \frac{1}{b - a} \int_{a}^{b} f(x) \mathrm{d}x.$\par
% p052-b03; source PDF page 52
\par\Needspace{5\baselineskip}
\blockanchor{p052-b03}
\majorquestion{五}解答题. 每题 15 分, 共 45 分.\par
% p052-b04; source PDF page 52
\blockanchor{p052-b04}
\question{1}讨论反常积分 $\int_0^{+\infty} \frac{\sin \sqrt{x}}{x^p} \, \mathrm{d}x$ 的敛散性, 其中参数 $p > \frac{1}{2}$.\par
% p052-b05; source PDF page 52
\blockanchor{p052-b05}
\question{2}讨论函数列  $f_{n}(x)=\frac{nx}{nx+1}, n=1,2,\cdots, x\in[0,+\infty)$  的一致收敛性及极限函数的连续性和可微性.\par
% p052-b06; source PDF page 52
\blockanchor{p052-b06}
\question{3}讨论级数 $\sum_{n=0}^{\infty}\left|\sin (\pi \sqrt{n^2 + 1})\right|^p$ 的敛散性, 其中 $p > 0$.\par
% p052-b07; source PDF page 52
\par\Needspace{5\baselineskip}
\blockanchor{p052-b07}
\majorquestion{六}(15 分) 取 $u = \frac{x}{x^2 + y^2}, v = \frac{y}{x^2 + y^2}$ 为新的自变量, 变换方程 $\frac{\partial^2 z}{\partial x^2} + \frac{\partial^2 z}{\partial y^2} = 0$, 其中所涉及的各阶导数均存在且连续.\par
